%%% FIELD no-csi-narrowed-statement
When \(\mathcal C=\varnothing\), put \(k_0=\prod_{i\in W}\mathcal X_i\). Then \(\mathcal K=\{k_0\}\), the empty sequence \(()\) is the only authorized trace, \(H_{k_0,()}=G\), and \(Q[D]_{k_0,()}\) is the ordinary district \(Q\)-factor computed in \(G\) for every relevant district \(D\). For every ordinary Shpitser--Pearl \(\mathrm{ID}\) run that succeeds, decorating each of its district-level \(Q\)-factor derivations by \((k_0,())\) gives SUCCESS derivations whose unchanged arithmetic assembly is the same observed-cell formula; erasure of these decorations recovers the original derivations. Moreover, the map \((F,F')\mapsto(F',F)\) is a bijection from the rooted C-forest obstruction predicate to the ordinary hedge predicate. Consequently every ordinary \(\mathrm{ID}\) success formula and every ordinary \(\mathrm{ID}\) failure hedge has, respectively, a SUCCESS encoding and an obstruction encoding. This is a syntactic conservativity statement; it does not assert that graphical \(\mathrm{ID}\) is semantically complete for every fixed-cardinality positive submodel.
%%% FIELD no-csi-replacement-statement
When \(\mathcal C=\varnothing\), put \(k_0=\prod_{i\in W}\mathcal X_i\). Then \(\mathcal K=\{k_0\}\), the empty sequence \(()\) is the only authorized trace, \(H_{k_0,()}=G\), and \(Q[D]_{k_0,()}\) is the ordinary district \(Q\)-factor computed in \(G\) for every relevant district \(D\). For every ordinary Shpitser--Pearl \(\mathrm{ID}\) run that succeeds, decorating each of its district-level \(Q\)-factor derivations by \((k_0,())\) gives SUCCESS derivations whose unchanged arithmetic assembly is the same observed-cell formula; erasure of these decorations recovers the original derivations. Moreover, the map \((F,F')\mapsto(F',F)\) is a bijection from the rooted C-forest obstruction predicate to the ordinary hedge predicate. Consequently every ordinary \(\mathrm{ID}\) success formula and every ordinary \(\mathrm{ID}\) failure hedge has, respectively, a SUCCESS encoding and an obstruction encoding. This is a syntactic conservativity statement; it does not assert that graphical \(\mathrm{ID}\) is semantically complete for every fixed-cardinality positive submodel.
%%% FIELD no-csi-proof
Let
\[
  \Omega_W:=\prod_{i\in W}\mathcal X_i .
\]
Each state space is nonempty, as required for a normalized probability table; hence \(\Omega_W\neq\varnothing\), including the case \(W=\varnothing\), when \(\Omega_W=\{()\}\).  By \(\texttt{def:activation-partition}\), the joint activation map for an empty CSI list is
\[
  T_{\varnothing}:\Omega_W\longrightarrow \{0,1\}^{\varnothing}=\{()\},
  \qquad T_{\varnothing}(w)=().
\]
Its sole nonempty fiber is
\[
  T_{\varnothing}^{-1}(())=\Omega_W.
\]
Thus, with \(k_0:=\Omega_W\),
\[
  \mathcal K=\{k_0\}.
\]

Now let \(t=(e_1,\ldots,e_\ell)\) be an authorized trace on \(k_0\).  By \(\texttt{def:authorized-trace}\), every \(e_j\) records an authorizing relation belonging to \(\mathcal C\).  Since \(\mathcal C=\varnothing\), no such \(e_j\) exists, so \(\ell=0\) and \(t=()\).  Conversely, the definition explicitly permits the empty sequence.  Therefore
\[
  \{t:t\text{ is authorized on }k_0\}=\{()\}.
\]
By \(\texttt{def:context-graph}\), starting from \(G\) and performing the empty list of deletions gives
\[
  H_{k_0,()}=G\setminus\varnothing=G.
\]
The cell \(k_0=\Omega_W\) imposes no restriction on a row argument, and the empty trace imposes no deletion or substitution.  Hence \(\texttt{def:context-q-factor}\) gives, coordinatewise,
\[
  Q[D]_{k_0,()}
  =
  \sum_{u_D}
  \prod_{a\in D\cup U_D}
  \theta_a\!\left(a_a\mid a_{\operatorname{pa}_G(a)}\right)
  =:Q_G[D],
\]
where \(U_D\) is the finite set of latent roots entering \(D\), and the sum ranges over their finite state spaces.  This is exactly the ordinary district \(Q\)-factor.

We next compare derivation syntax.  By \(\texttt{def:success-derivation}\), a CSI substitution node is legal only when its CSI identifier occurs in \(t\).  Since \(t=()\), no CSI substitution node is legal.  The remaining constructors are observed-cell leaves, finite marginalization, conditioning with a positive-denominator judgment, district restriction, and leaf removal.  These are precisely the constructors used by the district-level \(Q\)-factor recursion in \(\texttt{lem:shpitser-pearl-id-output}\).  Define a decoration map \(\mathsf d\) recursively: attach \((k_0,())\) to every node of an ordinary district-level derivation and leave its arithmetic operation and observed-cell index unchanged.  Define \(\mathsf e\) by erasing those two labels.  Structural induction on a derivation tree gives
\[
  \mathsf e(\mathsf d(\delta))=\delta
\]
at a leaf, and the same equality after each of marginalization, conditioning, district restriction, and leaf removal because \(\mathsf d\) changes none of their arguments.  Thus it holds for every ordinary district-level derivation.  Conversely, induction on the decorated image gives
\[
  \mathsf d(\mathsf e(\widetilde\delta))=\widetilde\delta.
\]
Therefore decoration is a bijection between ordinary district-level derivations and its decorated SUCCESS image, and arithmetic assembly across districts is unchanged.

The denominator judgments required by the SUCCESS grammar are available on the positive model class.  Indeed, by \(\texttt{ass:strict-cpt-positivity}\), for every complete observed assignment \(v\),
\[
  P_\theta(V=v)
  =\sum_{u\in\prod_{j\in U}\mathcal X_j}
    \prod_{a\in V\cup U}
      \theta_a\!\left(a_a\mid a_{\operatorname{pa}_G(a)}\right)>0.
\]
The index set is finite and nonempty, and every summand is a product of strictly positive coordinates.  Every marginal used as a conditioning denominator is a nonempty finite sum of such cells and is therefore strictly positive.  Hence every decorated successful formula is total in \(\mathscr E(P_{\mathrm{obs}})\); no boundary extension or limiting argument is used.

Finally consider \(\texttt{def:rooted-forest-obstruction}\).  Because \(H_{k_0,()}=G\), an obstruction pair satisfies
\[
  F\subseteq F'\subseteq G,
  \qquad F\cap X=\varnothing,
  \qquad F'\cap X\neq\varnothing,
\]
and both members are one-district C-forests with the same childless root set \(R_0\), where
\[
  R_0\subseteq\operatorname{An}_{G\setminus X}(Y).
\]
Set \(F_{\mathrm{large}}:=F'\) and \(F_{\mathrm{small}}:=F\).  Then
\[
  F_{\mathrm{small}}\subseteq F_{\mathrm{large}},
  \qquad
  F_{\mathrm{large}}\cap X\neq\varnothing,
  \qquad
  F_{\mathrm{small}}\cap X=\varnothing,
\]
with the same C-forest and root conditions.  These are the ordinary hedge conditions recorded in \(\texttt{lem:shpitser-pearl-id-output}\).  Swapping the two members again recovers \((F,F')\), so the map is a bijection.  The cited result says that an ordinary successful \(\mathrm{ID}\) run returns its observed-law formula and an ordinary failed run returns such a hedge.  Applying the decoration map in the first case and the forest-pair swap in the second proves the two claimed encodings.  The causal target remains the independently defined quantity \(q(\theta)=P_{\theta,\mathbf x}(\mathbf y)\); on a successful ordinary run, equality of the returned observed-law expression with this target follows from the cited soundness theorem and is not imposed by definition.  No rank, asymptotic, or limit condition is involved.
